%% Research Note of the AAS: blend ceiling for light-curve classifiers of long-period TESS candidates. 2026-11-25.
%% Text from DRAFT_RNAAS.md (commit 800412e), converted to AASTeX by hand; references checked on ADS 2026-11-25.
%% Source, numbers and provenance: https://github.com/coracle-us/duo-transits (rnaas/blend-ceiling/; see numbers.md).
\documentclass{aastex701}

\begin{document}

\title{Catalogued Eclipsing Binaries Flatter Light-curve Classifiers of Long-period TESS Candidates}

\author[orcid=0009-0000-7524-4470]{Jon McBee}
\affiliation{Independent researcher, USA}
\email{jon.mcbee@gmail.com}

\begin{abstract}
Classifiers of TESS candidates are often tested against catalogued eclipsing binaries (EBs). At short periods that
flatters them, because many real false positives are blends. We repeat the test at long periods ($P>10$\,d or single
transits): 336 confirmed or known planets against 215 TFOP false positives or false alarms. A gradient-boosting
classifier, never trained on TFOP false positives, reaches AUC 0.959 (0.945--0.971) against catalogued EBs and 0.959
against windows with no transit, but 0.649 (0.603--0.695) against TFOP false positives. Against those whose TOI
comments name a blend it scores 0.525 (0.441--0.607), consistent with chance. Comments flag 22\% of the false
positives as blends, which caps any light-curve-only classifier at AUC 0.89. Within depth bins the AUC is 0.57--0.63.
A first model without no-transit training windows failed a no-transit check but gave the same picture.
\end{abstract}

\keywords{\uat{Exoplanet detection methods}{489} --- \uat{Transit photometry}{1709} --- \uat{Eclipsing binary stars}{444}}

\section{Introduction}

Machine-learning classifiers for TESS candidates are usually scored against catalogued EBs.
\citet{PagesNavarrete2026} trained a CNN to detect ultra-short-period planets. It reached AUC 0.97 against
catalogued EBs but 0.50 against false positives rejected by the TESS Follow-up Observing Program (TFOP). At least
42\% of the false positives they examined were blends, which sets a heuristic ceiling $\mathrm{AUC_{max}}=1-f/2$ on
any classifier that sees only the target's light curve.

\citet{Salinas2026} introduced a CNN that estimates the depth, duration and mid-time of a single transit, with
calibrated uncertainties. They evaluated it on 120 confirmed planets with $P>10$\,d, but on no false positives, and
it gives no detection score, so we cannot score candidates with it. Here we run the two-negative-set test at long
periods with our own light-curve-only baseline, and measure the blend fraction among long-period TFOP false
positives.

\section{Data and method}

\textbf{Sample.} From the ExoFOP TOI table \citep{Guerrero2021}, downloaded 2026 November 24, we take TOIs with TFOP
Working Group disposition CP or KP (positives) or FP or FA (negatives), and $P>10$\,d or no period. Light curves are
SPOC 2-min \citep{Jenkins2016}, else QLP \citep{Huang2020}, one per sector, via Lightkurve
\citep{Lightkurve2018}. Windows are 99.5\,h long, binned to 30\,min and centred on each predicted transit with
timing error $3\sigma<25$\,h, and need at least half their bins filled. 336 positives and 215 negatives have at
least one valid window (median 8 per TOI).

\textbf{Training.} Positives are the CP/KP windows. Negatives are 410 EBs from \citet{Prsa2022} with $P>10$\,d that
are not TOIs, at their primary eclipses; 1146 ``TCE-like'' windows, the highest-SNR box dip in each light curve of
those stars at least 3\,d from any known transit or eclipse; and 1182 ``quiet'' windows, random windows from the
same stars under the same 3\,d rule. TFOP false positives are never used in training.

\textbf{Classifier.} Gradient boosting \citep{Pedregosa2011} on 15 features from the window flux alone: box-search
SNR and depth, duration, noise and red-noise ratio, V/U shape, asymmetry, second-dip SNR, skewness, brightening and
gap fraction. It is cross-fitted over 5 folds by TIC, so no star is scored by a model that saw it. An object's score
is its maximum over its windows.

\textbf{Blends.} A false positive is a blend under rule C if its TOI comments name a nearby EB, a centroid offset or
an off-target source, and under rule G if a Gaia DR3 \citep{Gaia2023} neighbour within $42''$ could produce the TOI
depth on its own. We give $f$ and the ceiling with Wilson intervals, and AUCs with 2000-sample bootstrap 95\% CIs.

\textbf{Order of work.} The rules were committed before any window was scored. The first model (A) had no quiet
windows. It failed the pre-registered no-transit check: 50 random out-of-transit windows from planet hosts scored a
median of 0.91, against 0.99 for planets, because every training negative was a strong dip. We then added the quiet
class as a dated rule change, with a pass criterion fixed in advance (median no-transit score below 0.5 and below the
false-positive median), and retrained once (model B). We report B, with A in brackets.

\section{Results}

\textbf{Blend fraction.} Under rule C, 47 of the 216 false positives are blends (215 have a valid window),
$f=0.218$ (0.168--0.277), so $\mathrm{AUC_{max}}=0.891$ (0.861--0.916): about half the fraction at ultra-short
periods. Rule G alone flags 70\% of the false positives but also 60\% of the confirmed planets, so it barely tells
them apart; C or G gives $f=0.718$ and $\mathrm{AUC_{max}}=0.641$, which we regard as an overcount.

\textbf{AUCs}, positives against each negative set (Figure~\ref{fig:auc}): catalogued EBs 0.959 (0.945--0.971)
[0.949]; no-transit windows 0.959 (0.939--0.976) [0.786]; all TFOP false positives 0.649 (0.603--0.695) [0.643];
non-blend (rule C) 0.684 (0.637--0.730) [0.668]; blend (rule C) 0.525 (0.441--0.607) [0.555].

\textbf{Checks.} Model B passes the no-transit check: the 50 windows score a median of 0.46 (10--90\%: 0.21--0.74),
against 0.98 for planets and 0.97 for TFOP false positives. A feature model without the three depth-scale features
and a small 1D CNN, trained the same way, fail it (medians 0.56 and 0.68); they reach 0.664 and 0.629 against all
false positives. The false positives are deeper than the planets (median 4090 vs 1680\,ppm, KS $p=2\times10^{-9}$)
and slightly fainter (Tmag 10.4 vs 10.1). Within four fixed depth bins the AUC is 0.57--0.63, with lower 95\% limits
of 0.49--0.52, so much of the 0.649 comes from depth.

\section{Discussion}

At long periods, as at short ones, catalogued EBs flatter a light-curve classifier: AUC 0.96 against them, 0.65
against real TFOP false positives. Blends are scored at chance. Teaching the model what an empty window looks like
fixed the no-transit check (0.79 to 0.96) but left the TFOP comparison where it was: easy negatives say little about
the false positives TFOP actually rejects.

The blend ceiling from TOI comments, 0.89, sits well above what this baseline reaches, so for it blends are not yet
the binding limit; the non-blend false positives, scored at 0.68, are. Any long-period classifier, including a
detection head added to networks like that of \citet{Salinas2026}, should be scored against TFOP false positives,
within depth bins, and on no-transit windows. Sample, rules, code and scores are at
\url{https://github.com/coracle-us/duo-transits} (directory \texttt{rnaas/blend-ceiling}).

\begin{figure}[ht!]
\plotone{fig1.pdf}
\caption{Left: ROC curves of the positives against all TFOP false positives, non-blend and blend false positives
(rule C), for model B. Right: AUC with 95\% bootstrap CIs per negative set, for model B (filled) and model A (open).
Horizontal lines are the blend ceilings $1-f/2$ (bands: Wilson intervals) for rule C and for C or G; the dashed line
is chance.\label{fig:auc}}
\end{figure}

\begin{acknowledgments}
This work uses TESS SPOC \dataset[doi:10.17909/t9-nmc8-f686]{https://doi.org/10.17909/t9-nmc8-f686} and QLP
\dataset[doi:10.17909/t9-r086-e880]{https://doi.org/10.17909/t9-r086-e880} light curves from MAST, ExoFOP
\dataset[doi:10.26134/ExoFOP5]{https://doi.org/10.26134/ExoFOP5}, and Gaia DR3 served by the ARI Heidelberg TAP
mirror.
\end{acknowledgments}

\facilities{TESS, Gaia}

\software{Lightkurve \citep{Lightkurve2018}, scikit-learn \citep{Pedregosa2011}, astroquery, NumPy, SciPy, pandas,
Matplotlib, PyTorch}

\section*{AI disclosure statement}
I used Anthropic's Claude, an AI assistant, under my supervision to write and run the analysis code and to draft and
edit this note. I reviewed the text, checked every number against the saved outputs and every reference against ADS,
and take full responsibility for the content. The figure was made by the analysis code from the data.

\begin{thebibliography}{}
\bibitem[Gaia Collaboration et al.(2023)]{Gaia2023} Gaia Collaboration, Vallenari, A., Brown, A.~G.~A., et al. 2023, A\&A, 674, A1
\bibitem[Guerrero et al.(2021)]{Guerrero2021} Guerrero, N.~M., Seager, S., Huang, C.~X., et al. 2021, ApJS, 254, 39
\bibitem[Huang et al.(2020)]{Huang2020} Huang, C.~X., Vanderburg, A., P{\'a}l, A., et al. 2020, RNAAS, 4, 204
\bibitem[Jenkins et al.(2016)]{Jenkins2016} Jenkins, J.~M., Twicken, J.~D., McCauliff, S., et al. 2016, Proc. SPIE, 9913, 99133E
\bibitem[Lightkurve Collaboration et al.(2018)]{Lightkurve2018} Lightkurve Collaboration, Cardoso, J.~V.~d.~M., Hedges, C., et al. 2018, ascl:1812.013
\bibitem[Pag{\'e}s Navarrete(2026)]{PagesNavarrete2026} Pag{\'e}s Navarrete, J.~F. 2026, arXiv:2610.42009
\bibitem[Pedregosa et al.(2011)]{Pedregosa2011} Pedregosa, F., Varoquaux, G., Gramfort, A., et al. 2011, JMLR, 12, 2825
\bibitem[Pr{\v s}a et al.(2022)]{Prsa2022} Pr{\v s}a, A., Kochoska, A., Conroy, K.~E., et al. 2022, ApJS, 258, 16
\bibitem[Salinas et al.(2026)]{Salinas2026} Salinas, H., Brahm, R., Olmschenk, G., et al. 2026, arXiv:2610.24882
\end{thebibliography}

\end{document}
